\documentclass[11pt]{article}
\usepackage[a4paper,margin=1in]{geometry}
\usepackage{amsmath,amsthm,amssymb}
\usepackage{microtype}
\usepackage{hyperref}

\title{Introduction and Conclusion \\
       \large{[draft for review, with venue notes]}}
\author{}
\date{}

\begin{document}
\maketitle

\noindent\textbf{Working title.}
\emph{Step-Indexed Hoare Logic as Ramified Forcing.}
Alternative: \emph{A Mechanized Ramified-Forcing Account of Step-Indexed Hoare Logic} --- includes ``mechanized'' for ITP/CPP audiences but is wordier. I default to the shorter form below.

\section{Introduction}

Step-indexed reasoning has become the standard semantic technique
for verifying higher-order programming languages with mutable state,
recursive types, and concurrency. From Appel and McAllester's
original step-indexed model of recursive types to the modern Iris
ecosystem, the same basic idea recurs: instead of asking whether a
program satisfies a specification, ask whether it satisfies the
specification \emph{at every finite observation depth $n$}, then
recover the unindexed statement as the limit. Each level of the
hierarchy of approximations is defined without circularity by appeal
only to lower levels, and a ``later'' modality $\triangleright$
mediates the descent from one level to the next. The technique is
sound, comprehensive, and widely mechanized; in Iris alone it
underwrites verifications of substantial concurrent programs.

What the step-indexed literature does not usually say is that this
technique is, in classical set-theoretic terms, \emph{forcing}.
Cohen introduced forcing in 1963 to prove the independence of the
continuum hypothesis. In its original presentation --- called
\emph{ramified forcing} after the ramified type theory of Russell
and Whitehead --- the construction of the generic extension is
stratified by a well-founded hierarchy of \emph{levels}, with
recursion at level $\alpha$ permitted to consult only levels strictly
below $\alpha$. Shoenfield's 1971 reformulation eliminated the
ramification for the set-theoretic application: the modern textbook
presentation of forcing does not use it. But the structure ---
well-founded stratification, level descent, a forcing relation
$p \Vdash \varphi$ that is monotone in the strength of $p$ --- is
exactly the structure of step-indexed logic.

In this paper we make the correspondence precise. The forcing
\emph{conditions} of step-indexed Hoare logic are the natural numbers
$\omega$ with the order $\le$; the forcing \emph{relation}
$n \Vdash \varphi$ is monotone in $n$ (downward closure); the
\emph{later} modality $\triangleright$ is the level descent operator
of ramified forcing; and Löb's rule, which underwrites all
recursive reasoning in step-indexed systems, is precisely
well-founded induction on the forcing order, internalized into the
logic.

\paragraph{Why does this matter?}
The forcing reading is not just a re-labelling. It clarifies three
things that are otherwise scattered across the step-indexed
literature.

First, it explains \emph{why} the technique has the shape it does.
The well-founded stratification is not a presentation convenience
to be eliminated; it is constitutive of the soundness of the modal
logic. Section~\ref{sec:lob-essential} makes this precise: we
exhibit a non-well-founded forcing notion (the integers $\mathbb{Z}$)
where every other piece of the framework goes through, but Löb's
rule fails.

Second, it diagnoses where the internal Löb rule earns its keep.
Section~\ref{sec:imp-vs-lam} compares two mechanized soundness
proofs of recursive Hoare rules: the \texttt{while} rule for IMP
(Section~\ref{sec:imp}) and the \texttt{fix} rule for the untyped
$\lambda$-calculus (Section~\ref{sec:lam}). The first admits an
unproblematic meta-level proof by induction on the step index; the
second does not, and the natural proof invokes the internal Löb rule
directly. The structural reason --- that the soundness predicate is
itself recursive in the \texttt{fix} case --- is made precise by
the forcing reading.

Third, it identifies step-indexed Hoare logic as the propositional
fragment of an existing construction in the type-theory literature:
the Jaber--Tabareau--Sozeau forcing translation of CIC.
Section~\ref{sec:bridge} formalizes this identification via a
constant embedding $\mathsf{Prop} \hookrightarrow \mathtt{iProp}$
that is a Heyting algebra homomorphism, with $\triangleright$
appearing as the step-indexed-specific addition.

\paragraph{Contributions.}

\begin{enumerate}

\item A mechanized re-presentation of step-indexed Hoare logic
as ramified forcing, in Rocq.
The propositional layer (Section~\ref{sec:prop}) gives a forcing
semantics for an intuitionistic logic with the later modality and
Löb's rule. The Hoare layer is constructed over IMP
(Section~\ref{sec:imp}) and over an untyped $\lambda$-calculus with
general recursion (Section~\ref{sec:lam}).

\item An explicit comparison of the IMP \texttt{while} rule and the
$\lambda$-calculus \texttt{fix} rule that diagnoses where the
internal Löb rule is essential
(Section~\ref{sec:imp-vs-lam}). For IMP, soundness goes through
either via meta-level induction on the step index or via the
internal Löb rule; we exhibit both proofs and observe their
structural identity. For the $\lambda$-calculus with \texttt{fix},
the natural soundness proof must invoke the internal Löb rule,
because the soundness predicate is itself recursive in the
program. This is the strongest piece of evidence for the slogan that
``Löb is what makes step-indexed Hoare logic work in higher-order
settings.''

\item An abstract forcing framework
(Section~\ref{sec:abstract}) parametric over an arbitrary
well-founded forcing structure. We exhibit two instances:
$(\omega, \le)$, recovering the standard step-indexed semantics; and
$(\omega \times \omega, \le_{\mathrm{pw}})$, with the pointwise
order, modelling independent step-indices for parallel components.
The internal Löb rule, in the abstract setting, is proved by
\texttt{well\_founded\_ind} on the strict refinement order. This
makes precise the equivalence: \emph{internal Löb is well-founded
induction on the forcing order, internalized}.

\item A small new metatheoretic result: \emph{well-foundedness of
the forcing notion is essential}
(Section~\ref{sec:lob-essential}). We exhibit a non-well-founded
``would-be'' forcing structure over $\mathbb{Z}$, in which every
iProp construction goes through except Löb's rule. The premise of
Löb's rule $(\triangleright \bot) \to \bot$ is forced at every
condition --- vacuously, since $\mathbb{Z}$ has no minimum --- while
$\bot$ is forced at none. This is a special case of the model theory
of the Gödel--Löb provability logic \textsf{GL}, brought into the
step-indexed Hoare logic setting and mechanized.

\item A partial bridge to the Jaber--Tabareau--Sozeau forcing
translation of CIC (Section~\ref{sec:bridge}). We formalize the
constant embedding $\mathsf{Prop} \hookrightarrow \mathtt{iProp}$
as a Heyting algebra homomorphism, and observe that the $\triangleright$
modality is the genuinely new content of the iProp universe over
the embedded image of Prop. A full mechanization of the type-level
JTS translation is future work.

\end{enumerate}

\paragraph{What is and is not new.}
The individual ingredients are largely known to specialists in
neighbouring fields. The topos-of-trees perspective on step-indexing
(Birkedal et al.) is implicitly a forcing semantics; the Gödel--Löb
modal logic of provability (Solovay, Smoryński, Boolos) is the modal
logic of provability and well-foundedness; the JTS forcing
translation of CIC is the type-theoretic analog of forcing. What is
new is the assembly: a single mechanized account that makes the
forcing reading explicit, diagnoses the role of the internal Löb
rule, and gives a parametric framework with the metatheoretic
guarantee that well-foundedness cannot be relaxed.

\paragraph{Outline.}
Section~\ref{sec:prop} gives the propositional layer of the forcing
semantics, including the Heyting-algebra connectives, the later
modality, and Löb's rule. Section~\ref{sec:imp} develops Hoare
logic for IMP and proves the \texttt{while} rule by meta-level
induction; Section~\ref{sec:lam} develops Hoare logic for an untyped
$\lambda$-calculus with \texttt{fix} and proves the soundness rule
for \texttt{fix} via the internal Löb rule. Section~\ref{sec:imp-vs-lam}
compares these two proofs explicitly and articulates the diagnostic
claim. Section~\ref{sec:abstract} gives the abstract forcing
framework. Section~\ref{sec:lob-essential} proves Löb's rule fails
without well-foundedness. Section~\ref{sec:bridge} establishes the
constant-embedding bridge to JTS forcing translations.
Section~\ref{sec:related} positions the contribution against the
literature. Section~\ref{sec:conclusion} concludes.

\paragraph{Rocq development.}
The complete Rocq development is approximately 800 lines across 11
files, building cleanly against the Rocq standard library only (no
external libraries such as Iris or stdpp). Per-file inventory is in
the project README.

\section{Conclusion}

We have given a forcing-style semantics for step-indexed Hoare
logic, mechanized in Rocq across the propositional layer, an IMP
Hoare logic, and an untyped $\lambda$-calculus with general
recursion. The forcing reading places the step-indexed technique
in direct correspondence with Cohen's ramified forcing construction:
the natural numbers are the forcing conditions, the later modality
$\triangleright$ is the level-descent operator, and Löb's rule is
well-founded induction on the forcing order, internalized into the
logic.

\paragraph{What we found.}
Two observations stand out. First, in the first-order setting of
IMP, the \texttt{while} rule admits two structurally identical
soundness proofs: an outer induction on the step index, and an
appeal to the internal Löb rule. The two proofs have the same
content because at the meta level, Löb's rule \emph{is} induction
on the step index. Mechanizing both side by side, as we do, makes
the equivalence concrete rather than analogical.

Second, in the higher-order setting of the untyped $\lambda$-calculus
with \texttt{fix}, the analogous soundness proof has a different
shape. The soundness predicate for a recursive function is itself
defined in terms of the function being verified, and unrolling this
self-reference at the meta level produces a non-trivial induction
that does not arise from the program's step count alone. The natural
proof, which we mechanize, invokes the internal Löb rule directly:
``show that under the assumption $\triangleright \varphi$ we have
$\varphi$, then conclude $\varphi$ by Löb.'' No meta-level induction
on a step count is needed. This is, in our opinion, the simplest
diagnostic for the slogan that \emph{Löb is what makes step-indexed
Hoare logic work for higher-order recursion}.

\paragraph{Three further observations.}
The abstract framework of Section~\ref{sec:abstract} makes precise
that the choice of $\omega$ as the step-index domain is a
\emph{computational} convenience rather than a foundational one. Any
well-founded forcing notion supports the same internal logic, with
$\triangleright$ defined uniformly as ``forces at every strict
predecessor'' and Löb derived by well-founded induction. We exhibit
the pointwise product $\omega \times \omega$ as a second instance, but
the framework is generic.

Section~\ref{sec:lob-essential} shows that the well-foundedness
cannot be dropped: we construct a non-well-founded forcing notion
over $\mathbb{Z}$ in which Löb's rule fails. This is a special case
of the model theory of \textsf{GL} (Smoryński, Boolos), but the
statement for step-indexed Hoare logic does not, as far as we are
aware, appear in the verification literature. The claim has
methodological weight: step-indexed Hoare logic is not just
``Hoare logic with a step counter''; it requires a well-founded
forcing notion to be sound at all.

Section~\ref{sec:bridge} verifies that the iProp framework is the
propositional fragment of the Jaber--Tabareau--Sozeau forcing
translation of CIC, via a constant embedding
$\mathsf{Prop} \hookrightarrow \mathtt{iProp}$ that we prove to be
a Heyting algebra homomorphism. The $\triangleright$ modality is
the genuinely new content of iProp over the embedded image, and
arises uniformly from the well-founded strict refinement on
conditions.

\paragraph{Limitations.}
The Hoare logic for the $\lambda$-calculus in
Section~\ref{sec:lam} is, at present, restricted to partial
correctness against a single specification predicate $S$ that
serves as both pre- and post-condition. The fix soundness rule
\texttt{wp\_fix} (Section~\ref{sec:lam}) accordingly handles only
this case. A more general rule with distinct pre/post-conditions, a
rule for fix at higher types, and an integration with separation
logic on a heap are all natural directions but lie outside the
present development.

The bridge to JTS forcing translations of Section~\ref{sec:bridge}
formalizes only the constant embedding. A full mechanization of the
type-level translation, against an explicit forcing structure and
with the propositional rules verified to match our \texttt{iProp},
is the natural next step. This would make precise (rather than
heuristic) the claim that our abstract framework is the
propositional fragment of the JTS translation.

The abstract framework of Section~\ref{sec:abstract} parametrizes
the propositional layer over arbitrary well-founded forcing
structures, but the IMP and $\lambda$-calculus developments of
Sections~\ref{sec:imp}--\ref{sec:lam} are presented at the
specific $(\omega, \le)$ instance for clarity. Lifting them to be
parametric in the forcing structure is straightforward and not
done here.

\paragraph{Future directions.}
Several extensions suggest themselves. Most concretely:

\begin{itemize}

\item \emph{Mechanizing the full JTS translation} (or a useful
fragment) and verifying that its propositional fragment coincides
definitionally with our \texttt{iProp} at the corresponding forcing
structure. This would convert the heuristic correspondence of
Section~\ref{sec:bridge} into a theorem.

\item \emph{Step-indexing over richer forcing structures} ---
transfinite ordinals (cf.\ Transfinite Iris, Spies et al.),
branching conditions for genuinely concurrent step counters, or
mixed structures combining several step-index dimensions. The
abstract framework already supports these in principle; the work
is to exhibit non-trivial verifications at the new instances.

\item \emph{Lifting the Hoare-logic developments to the abstract
framework}, so that IMP and the $\lambda$-calculus admit
verifications parametric in the forcing structure. This is
straightforward but bookkeeping-heavy; it would make the
parametric story of Section~\ref{sec:abstract} the primary
presentation rather than an afterthought.

\item \emph{An integration with separation logic} along the lines
of Iris, lifting the forcing reading from the propositional logic
to the resource-aware extension. This is a substantial undertaking
but would let the forcing perspective speak directly to
state-of-the-art verification.

\end{itemize}

The broader point we hope the paper makes is that step-indexed
reasoning, viewed as a forcing construction, sits at the
intersection of three research traditions that do not usually talk
to each other: set-theoretic forcing, modal logic of provability,
and step-indexed semantics for verification. Each tradition contains
piece-of-the-picture insights that the other two could profit from.
The mechanization that accompanies this paper is intended both as
evidence for the technical claims and as a starting point for
others wishing to explore the intersection further.

\end{document}
